\chapter{Infinite Series and Convergence}

There is something faintly outrageous about the idea of adding up infinitely many numbers. Addition, as you first learned it, is a finite act: you line up two numbers, or ten, or a million, and you combine them. The process \emph{stops}. But what could it possibly mean to write
\[
1 + \frac{1}{2} + \frac{1}{4} + \frac{1}{8} + \frac{1}{16} + \cdots
\]
and ask for ``the sum''? You can never reach the end. There is no last term to add. And yet -- this is the quiet miracle the chapter is built around -- some of these endless sums settle down to a perfectly definite, ordinary number, while others run off to infinity or thrash around forever without ever choosing a value. Our job is to figure out \emph{which is which}, and, more importantly, to understand \emph{why}.

This is not an idle game. Every time you evaluate $e^{0.3}$ on a calculator, or your computer reports $\sin(1)$, or a physicist models the gravitational pull of a ring of mass, an infinite series is doing the work underneath. The functions you have grown to love -- $e^x$, $\sin x$, $\ln(1+x)$ -- \emph{are} infinite series in disguise, and in the next chapter we will pry that disguise off. But none of it makes sense until we are crystal clear about what it means for an infinite sum to have a value at all. So we begin at the only honest place: the partial sums.

\section{What It Means to Add Infinitely Many Things}
\label{ch29:partial-sums}

The trick -- and it is the single most important idea in the chapter, the seed from which everything else grows -- is that we never actually perform an infinite addition. We \emph{can't}. Instead we perform a sequence of perfectly ordinary finite additions and watch where they are heading.

\begin{intuitionbox}
Imagine walking toward a wall. First you cover half the remaining distance, then half of what's left, then half again. After each step you have walked a finite, knowable distance. You will never literally arrive, but it is intuitively clear \emph{where} you are arriving: at the wall. The infinite sum is not the act of taking infinitely many steps. It is the \emph{destination} those steps approach.
\end{intuitionbox}

Let us make this precise. Given an infinite list of numbers $a_1, a_2, a_3, \dots$ -- which we call the \emph{terms} -- we build a brand-new list out of them, the list of \textbf{partial sums}:
\[
S_1 = a_1, \qquad
S_2 = a_1 + a_2, \qquad
S_3 = a_1 + a_2 + a_3, \qquad \dots, \qquad
S_n = \sum_{k=1}^{n} a_k .
\]
Each $S_n$ is a finite sum -- no mystery, no infinity, just arithmetic. The infinite series is then \emph{defined} to be the limit of this sequence of partial sums, if that limit exists.

\begin{definition}[Convergence of a series]
\label{ch29:def-convergence}
Let $a_1, a_2, a_3, \dots$ be a sequence of real numbers and let $S_n = \sum_{k=1}^n a_k$ be its $n$th partial sum. We say the series $\sum_{k=1}^{\infty} a_k$ \textbf{converges} to the number $S$ if
\[
\lim_{n \to \infty} S_n = S,
\]
and we then write $\sum_{k=1}^{\infty} a_k = S$. If the sequence $(S_n)$ has no finite limit, we say the series \textbf{diverges}.
\end{definition}

\begin{happeningbox}
The symbol $\sum_{k=1}^{\infty} a_k$ looks like an instruction to do something infinitely many times. It is not. It is shorthand for a \emph{limit of a sequence}: the sequence of finite partial sums. When we write $\sum_{k=1}^\infty \tfrac{1}{2^k} = 1$, we are not claiming that somebody added up infinitely many fractions and got $1$. We are claiming something about a sequence of ordinary numbers $\tfrac12, \tfrac34, \tfrac78, \tfrac{15}{16}, \dots$ -- namely that they can be made as close to $1$ as you please by going far enough out. The infinite sum is a \emph{statement about approach}, dressed up to look like an act of addition. Once you internalize this, the paradoxes of infinity dissolve: there is no completed infinite process anywhere, only a limit.
\end{happeningbox}

So ``convergence'' and ``divergence'' are not really new ideas. They are the limit concept from earlier in the book, applied to the particular sequence $(S_n)$. Everything we prove about series is, underneath, a fact about the convergence of these partial-sum sequences. Keep that anchor in mind whenever a test starts to feel like a magic incantation.

\begin{figure}[H]
\centering
\begin{tikzpicture}
\begin{axis}[
    width=12cm, height=6cm,
    xlabel={$n$ (number of terms added)}, ylabel={partial sum $S_n$},
    xmin=0, xmax=10.5, ymin=0, ymax=2.2,
    xtick={1,2,3,4,5,6,7,8,9,10},
    samples=10, domain=1:10,
    legend pos=south east, legend cell align=left,
    grid=both, grid style={gray!18},
]
\addplot[only marks, mark=*, mark size=1.6pt, curveblue]
    coordinates {(1,1)(2,1.5)(3,1.75)(4,1.875)(5,1.9375)(6,1.96875)(7,1.984375)(8,1.9921875)(9,1.99609375)(10,1.998046875)};
\addlegendentry{$S_n = \sum_{k=0}^{n-1} 1/2^k$}
\addplot[tangentred, dashed, thick, domain=0:10.5] {2};
\addlegendentry{limit $= 2$}
\end{axis}
\end{tikzpicture}
\caption{The partial sums of $1 + \tfrac12 + \tfrac14 + \cdots$ climb toward $2$, halving the gap each step but never crossing. The infinite sum \emph{is} that ceiling.}
\label{ch29:fig-partialsums}
\end{figure}

\begin{walkbox}
Before reading on, settle a question in your head on a walk: the terms $\tfrac{1}{2^k}$ get smaller and smaller and eventually become negligibly tiny. Is ``the terms shrink to zero'' enough, all by itself, to guarantee the sum is finite? Hold your answer. By the end of Section~\ref{ch29:harmonic} it may surprise you.
\end{walkbox}

\section{The Geometric Series: The One You Can Sum Exactly}
\label{ch29:geometric}

Almost every infinite sum is impossible to evaluate in closed form -- you can prove it converges, but you cannot write down a tidy formula for the total. There is one glorious exception that we can sum \emph{exactly}, and it happens to be the most important series in all of mathematics: the \textbf{geometric series}, in which each term is a fixed multiple $r$ of the one before.

\[
a + ar + ar^2 + ar^3 + \cdots = \sum_{k=0}^{\infty} a\, r^k .
\]

The number $r$ is the \emph{ratio}. Our walk-toward-the-wall example was the case $a=1$, $r=\tfrac12$. Let us find the total honestly, by computing the partial sums and taking their limit -- never by waving hands at infinity.

\begin{focusbox}
\textbf{The summing trick.} Let $S_n = a + ar + ar^2 + \cdots + ar^{n-1}$ (that's $n$ terms). Multiply the whole thing by $r$:
\[
r S_n = ar + ar^2 + \cdots + ar^{n-1} + ar^{n}.
\]
Now subtract. Look at what happens: almost everything cancels, because $rS_n$ is just $S_n$ shifted over by one slot.
\[
S_n - rS_n = a - ar^{n} \quad\Longrightarrow\quad (1-r)S_n = a\,(1 - r^{n}).
\]
Provided $r \neq 1$, divide:
\[
\boxed{\,S_n = a\,\frac{1 - r^{n}}{1 - r}\,}.
\]
\end{focusbox}

This is an \emph{exact} formula for the finite partial sum -- valuable on its own. Now we take the limit. Everything hinges on what $r^n$ does as $n \to \infty$. If $\abs{r} < 1$, then $r^n \to 0$ (repeatedly multiplying by something smaller than $1$ in size drives you to zero). If $\abs{r} > 1$, then $\abs{r^n} \to \infty$ and the partial sums blow up. The boundary cases $r = 1$ and $r = -1$ misbehave in their own ways.

\begin{theorem}[The geometric series]
\label{ch29:thm-geometric}
The geometric series $\sum_{k=0}^{\infty} a\, r^k$ converges if and only if $\abs{r} < 1$, and in that case
\[
\sum_{k=0}^{\infty} a\, r^{k} = \frac{a}{1 - r}.
\]
If $\abs{r} \geq 1$ (and $a \neq 0$) the series diverges.
\end{theorem}

\begin{proof}
From the boxed formula, $S_n = a\dfrac{1 - r^n}{1-r}$ for $r \neq 1$. If $\abs r < 1$ then $r^n \to 0$, so $S_n \to \dfrac{a(1-0)}{1-r} = \dfrac{a}{1-r}$. If $\abs r > 1$ then $\abs{r^n}\to\infty$, so $S_n$ has no finite limit. If $r = 1$, the series is $a + a + a + \cdots$ and $S_n = na \to \pm\infty$. If $r = -1$, the partial sums bounce $a, 0, a, 0, \dots$ and never settle. In every case with $\abs r \ge 1$ the series diverges.
\end{proof}

\begin{whybox}
Why does the $\abs{r}<1$ condition do \emph{all} the work? Because the geometric series carries its own fate in its terms. The ``tail'' that remains after $n$ terms is itself a geometric series, $ar^n + ar^{n+1}+\cdots = r^n \cdot \frac{a}{1-r}$, a shrunken copy of the original scaled by $r^n$. When $\abs r < 1$ that scaling factor crushes the leftover tail to nothing; the destination is in sight and the remaining journey shrinks geometrically. When $\abs r \ge 1$ the tail is as big as -- or bigger than -- the whole, so there is never any ``settling down.'' The self-similarity of the geometric series is exactly why it is the one we can pin down precisely, and it is why $\abs r$ versus $1$ is the only thing that matters.
\end{whybox}

\begin{example}[The walk to the wall, finished]
With $a=1$, $r=\tfrac12$: the sum is $\dfrac{1}{1 - \tfrac12} = 2$. Note our list started at $k=0$, so this is $1 + \tfrac12 + \tfrac14 + \cdots = 2$, matching Figure~\ref{ch29:fig-partialsums}.
\end{example}

\begin{example}[A repeating decimal is a geometric series]
What is $0.7777\ldots$ really? It is
\[
\frac{7}{10} + \frac{7}{100} + \frac{7}{1000} + \cdots = \sum_{k=1}^{\infty} \frac{7}{10}\left(\frac{1}{10}\right)^{k-1},
\]
a geometric series with first term $a = \tfrac{7}{10}$ and ratio $r = \tfrac{1}{10}$. So its value is
\[
\frac{a}{1-r} = \frac{7/10}{1 - 1/10} = \frac{7/10}{9/10} = \frac{7}{9}.
\]
Every repeating decimal is secretly a geometric series, which is why every repeating decimal is a fraction. The mysterious $0.999\ldots = 1$ is just $\dfrac{9/10}{1-1/10}=1$ -- not a trick, just the geometric formula.
\end{example}

\begin{historybox}
The geometric series is ancient. Around 250 BC, Archimedes summed $1 + \tfrac14 + \tfrac{1}{16} + \cdots = \tfrac43$ to find the area of a parabolic segment -- a genuine infinite sum, handled with total rigor, two thousand years before ``limit'' had a definition. He did not say ``let $n \to \infty$''; instead he proved by his \emph{method of exhaustion} that the area could be neither more nor less than $\tfrac43$, squeezing it from both sides. The intuition of convergence was available to the very best mathematicians long before the machinery. What the nineteenth century added was not the answer but the \emph{license}: a clean definition (our partial sums) that let ordinary mathematicians sum series safely, without an Archimedes' worth of cleverness each time.
\end{historybox}

\section{Telescoping Series: When the Middle Cancels}
\label{ch29:telescoping}

Besides the geometric series, there is one more family we can sum exactly, and it relies on the same essential move: arrange for the partial sum $S_n$ to collapse to just a few surviving terms. These are the \textbf{telescoping series}, named for the way an old collapsible spyglass folds down to almost nothing.

The idea: suppose each term can be written as a \emph{difference} of consecutive values of some sequence, $a_k = b_k - b_{k+1}$. Then
\[
S_n = \sum_{k=1}^n (b_k - b_{k+1}) = (b_1 - b_2) + (b_2 - b_3) + (b_3 - b_4) + \cdots + (b_n - b_{n+1}).
\]
Every interior term appears twice with opposite signs and cancels. Only the very first and very last survive:
\[
S_n = b_1 - b_{n+1}.
\]
The infinite sum is then $\lim_{n\to\infty} S_n = b_1 - \lim_{n\to\infty} b_{n+1}$, easy as soon as you know where $b_k$ heads.

\begin{example}[The textbook telescope]
Consider $\displaystyle\sum_{k=1}^{\infty} \frac{1}{k(k+1)}$. The terms shrink fast, but is there a sum? Use partial fractions -- the key that unlocks the difference structure:
\[
\frac{1}{k(k+1)} = \frac{1}{k} - \frac{1}{k+1}.
\]
So with $b_k = \tfrac1k$ we have $a_k = b_k - b_{k+1}$, and
\[
S_n = \left(1 - \tfrac12\right) + \left(\tfrac12 - \tfrac13\right) + \cdots + \left(\tfrac1n - \tfrac{1}{n+1}\right) = 1 - \frac{1}{n+1}.
\]
As $n\to\infty$, $\tfrac{1}{n+1}\to 0$, so the series sums to exactly $1$. Clean and complete.
\end{example}

\begin{trapbox}
Telescoping is seductive, and the trap is assuming \emph{everything} cancels except the first term. It does not. Always write out $S_n$ honestly and see which terms survive at \emph{both} ends. Some telescopes leave two survivors on each side (when the difference skips a step, like $b_k - b_{k+2}$), giving $S_n = b_1 + b_2 - b_{n+1} - b_{n+2}$. And the sum exists only if the trailing survivors have a limit. If $b_k$ does not converge, the telescope diverges no matter how prettily the middle cancels.
\end{trapbox}

\section{The First Thing to Check: The $n$th-Term Test}
\label{ch29:nth-term}

Before reaching for any heavy machinery, there is a one-second sanity check that can kill a series outright. It comes from a simple observation about \emph{convergent} series and their terms.

\begin{intuitionbox}
If the partial sums $S_n$ are settling down to a fixed number $S$, then consecutive partial sums must be getting closer and closer to each other. But the gap between consecutive partial sums is exactly one term: $S_n - S_{n-1} = a_n$. So the terms must be shrinking to zero. A convergent sum \emph{cannot} keep adding non-negligible chunks forever -- if it did, it would never settle.
\end{intuitionbox}

\begin{theorem}[The $n$th-term test for divergence]
\label{ch29:thm-nth-term}
If $\sum_{k=1}^{\infty} a_k$ converges, then $\lim_{n\to\infty} a_n = 0$. Equivalently (contrapositive): if $a_n$ does \emph{not} tend to $0$, the series diverges.
\end{theorem}

\begin{proof}
Suppose the series converges to $S$, so $S_n \to S$ and also $S_{n-1} \to S$ (same sequence, shifted). Then
\[
a_n = S_n - S_{n-1} \;\longrightarrow\; S - S = 0. \qedhere
\]
\end{proof}

\begin{trapbox}
This is the single most misused fact in the chapter, so read carefully. The theorem says one thing and one thing only: \emph{if the terms don't go to zero, the series diverges.} It does \textbf{not} say the converse. ``The terms go to zero'' does \textbf{not} imply the series converges. That converse is false, spectacularly so, and the harmonic series in Section~\ref{ch29:harmonic} is the famous counterexample. The $n$th-term test is a \emph{divergence detector only}. When the terms shrink to zero, the test is silent -- it tells you nothing, and you must reach for a real test.
\end{trapbox}

\begin{example}
$\sum \dfrac{n}{2n+1}$ diverges, because $\dfrac{n}{2n+1}\to \tfrac12 \neq 0$. We are adding chunks that stay near $\tfrac12$ forever; no settling is possible. Likewise $\sum (-1)^n$ diverges since the terms don't approach $0$ (they bounce between $1$ and $-1$). One glance, done.
\end{example}

\section{The Harmonic Series: A Genuine Surprise}
\label{ch29:harmonic}

Here is the series that should unsettle you, the one the walking question pointed toward:
\[
\sum_{n=1}^{\infty} \frac{1}{n} = 1 + \frac{1}{2} + \frac{1}{3} + \frac{1}{4} + \frac{1}{5} + \cdots
\]
The terms shrink to zero. They shrink, in fact, quite obediently. The $n$th-term test says nothing against it. Surely such a sum, built of ever-tinier pieces, must add up to something finite? It does not. The harmonic series \textbf{diverges} -- its partial sums crawl off to infinity, slowly but inexorably. This is one of the great surprises of mathematics, and the proof is so clever and so clean that it is worth knowing by heart.

\begin{focusbox}
\textbf{Oresme's grouping (c.~1350).} Group the terms in blocks whose lengths double, and bound each block from below:
\[
\underbrace{1}_{} + \underbrace{\frac12}_{} + \underbrace{\frac13 + \frac14}_{> \frac14 + \frac14 = \frac12} + \underbrace{\frac15 + \frac16 + \frac17 + \frac18}_{> \frac18\cdot 4 = \frac12} + \underbrace{\frac19 + \cdots + \frac1{16}}_{> \frac{1}{16}\cdot 8 = \frac12} + \cdots
\]
In each block, replace every term by the \emph{smallest} term in that block; the block still beats $\tfrac12$. So the sum contains infinitely many chunks each exceeding $\tfrac12$. Add enough blocks and you pass any number you like. The partial sums are unbounded; the series diverges.
\end{focusbox}

\begin{historybox}
Nicole Oresme, a French scholar of the fourteenth century, produced exactly this argument around 1350 -- centuries before calculus existed. It is one of the earliest truly rigorous results about an infinite process. The divergence is agonizingly slow: to push the partial sum past $10$ you need over twelve thousand terms; to reach $100$ you need more terms than there are atoms in your body. This sluggishness is precisely why the result is shocking. Everything \emph{looks} like it should converge. It is the slowest, gentlest possible march to infinity -- but it gets there.
\end{historybox}

\begin{aestheticbox}
The harmonic series teaches the deepest lesson of the whole subject: \emph{the rate at which terms shrink is everything.} The terms $\tfrac1n$ go to zero, but not fast enough; the terms $\tfrac{1}{2^n}$ go to zero fast enough. The entire theory of convergence -- every test that follows -- is really one question asked in many disguises: \emph{do the terms shrink fast enough?} Geometric decay ($r^n$) is fast enough. Decay like $\tfrac1n$ is not. The borderline lies somewhere in between, and mapping that borderline precisely is what the rest of the chapter is about.
\end{aestheticbox}

\section{The Integral Test: Series as Areas}
\label{ch29:integral-test}

The harmonic series begs a follow-up: where, exactly, is the borderline? We know $\sum \tfrac1n$ diverges. What about $\sum \tfrac{1}{n^2}$, or $\sum \tfrac{1}{n^{1.5}}$, or $\sum \tfrac{1}{\sqrt n}$? To settle a whole family at once we connect sums to integrals -- two ways of measuring the same underlying quantity, area under a decreasing curve.

\begin{intuitionbox}
Picture the terms $a_n = f(n)$ of a positive, decreasing sequence as the heights of rectangles, each of width $1$. The sum $\sum f(n)$ is the total area of an infinite staircase of rectangles. The integral $\int_1^\infty f(x)\,\dd x$ is the area under the smooth curve $y=f(x)$. Because the curve is decreasing, the rectangles and the curve trap each other: the staircase is never far from the smooth area. So the sum is finite exactly when the integral is finite. We've traded a hard sum for an integral we may know how to do.
\end{intuitionbox}

\begin{figure}[H]
\centering
\begin{tikzpicture}
\begin{axis}[
    width=12cm, height=6cm,
    xlabel={$x$}, ylabel={$y$},
    xmin=0.5, xmax=6.5, ymin=0, ymax=1.15,
    xtick={1,2,3,4,5,6}, ytick={0,0.5,1},
    samples=120, domain=1:6.2,
    legend pos=north east, legend cell align=left,
]
% rectangles: left-endpoint (overestimate), height f(n) on [n,n+1]
\addplot[ybar interval, fill=areafill, fill opacity=0.45, draw=curveblue!60]
    coordinates {(1,1)(2,0.5)(3,0.3333)(4,0.25)(5,0.2)(6,0.1667)};
\addlegendentry{rectangles $f(n)\cdot 1$}
\addplot[curveblue, very thick, domain=1:6.2] {1/x};
\addlegendentry{$y = 1/x$}
\end{axis}
\end{tikzpicture}
\caption{Each rectangle has height $f(n)=1/n$ and sits \emph{above} the curve on $[n,n+1]$. So $\sum_{n=1}^{N} f(n) > \int_1^{N+1} f(x)\,\dd x$. Since $\int_1^\infty \tfrac1x\,\dd x = \infty$, the harmonic sum is forced to infinity too -- a second proof of its divergence, by picture.}
\label{ch29:fig-integraltest}
\end{figure}

\begin{theorem}[The integral test]
\label{ch29:thm-integral}
Let $f$ be positive, continuous, and decreasing on $[1,\infty)$, and set $a_n = f(n)$. Then the series $\sum_{n=1}^{\infty} a_n$ and the improper integral $\int_1^{\infty} f(x)\,\dd x$ either both converge or both diverge.
\end{theorem}

\begin{proof}
Because $f$ decreases, on the interval $[n, n+1]$ we have $f(n+1) \le f(x) \le f(n)$. Integrating across that unit interval (width $1$) gives the rectangle-vs-curve sandwich:
\[
f(n+1) \;\le\; \int_{n}^{n+1} f(x)\,\dd x \;\le\; f(n).
\]
Sum the right inequality for $n=1$ to $N-1$: $\int_1^{N} f \le \sum_{n=1}^{N-1} f(n)$. Sum the left inequality for $n=1$ to $N-1$: $\sum_{n=2}^{N} f(n) \le \int_1^{N} f$. Together,
\[
\sum_{n=2}^{N} a_n \;\le\; \int_1^{N} f(x)\,\dd x \;\le\; \sum_{n=1}^{N-1} a_n .
\]
If the integral converges (finite as $N\to\infty$), the left inequality bounds the partial sums above; a bounded increasing sequence converges, so the series converges. If the integral diverges, the right inequality forces the partial sums to infinity, so the series diverges. The integral and the series are locked together.
\end{proof}

\begin{whybox}
The integral test works because a sum is a \emph{crude integral} and an integral is a \emph{refined sum} -- they are two computations of the same area, differing only by the gap between a staircase and the curve it hugs. When $f$ decreases, that gap is itself controlled (it is squeezed between consecutive terms), so it can never accumulate enough to make the sum and the integral disagree about being finite. This is the deep reason the test needs ``decreasing'': without it, the rectangles could wander far from the curve and the comparison would collapse. The hypothesis is not red tape; it is the very thing that keeps staircase and curve shackled together.
\end{whybox}

The payoff is an entire family settled in one stroke -- the single most useful convergence fact to memorize.

\begin{theorem}[The $p$-series]
\label{ch29:thm-pseries}
The series $\displaystyle\sum_{n=1}^{\infty} \frac{1}{n^p}$ converges if and only if $p > 1$.
\end{theorem}

\begin{proof}
Apply the integral test to $f(x)=x^{-p}$ (positive, decreasing for $x\ge 1$ when $p>0$; and for $p\le 0$ the terms don't even go to zero, so it diverges by the $n$th-term test). For $p \ne 1$,
\[
\int_1^{\infty} x^{-p}\,\dd x = \lim_{b\to\infty} \frac{x^{1-p}}{1-p}\Big|_1^{b} = \lim_{b\to\infty}\frac{b^{1-p}-1}{1-p}.
\]
When $p>1$, the exponent $1-p$ is negative, so $b^{1-p}\to 0$ and the integral equals $\tfrac{1}{p-1}$, finite: convergence. When $p<1$, $b^{1-p}\to\infty$: divergence. When $p=1$, $\int_1^\infty \tfrac1x\,\dd x = \lim_{b\to\infty}\ln b = \infty$: divergence (the harmonic series, again).
\end{proof}

\begin{focusbox}
\textbf{Memorize this.} $\sum \tfrac{1}{n^p}$ converges $\iff p > 1$. So $\sum\tfrac1{n^2}, \sum\tfrac1{n^3}, \sum\tfrac{1}{n^{1.01}}$ converge; $\sum\tfrac1n, \sum\tfrac{1}{\sqrt n}=\sum n^{-1/2}, \sum\tfrac{1}{n^{0.99}}$ diverge. The borderline sits exactly at $p=1$, and the harmonic series sits exactly on the borderline -- on the \emph{divergent} side. The borderline you have been hunting since Section~\ref{ch29:harmonic} is now located: it is $p=1$, and $1/n$ is the worst term that still diverges. Figure~\ref{ch29:fig-pseries} shows the three regimes side by side.
\end{focusbox}

\begin{figure}[H]
\centering
\begin{tikzpicture}
\begin{axis}[
    width=12cm, height=6cm,
    xlabel={$N$ (number of terms)}, ylabel={partial sum},
    xmin=0, xmax=20.5, ymin=0, ymax=3.7,
    xtick={2,5,10,15,20},
    legend pos=north west, legend cell align=left,
    grid=both, grid style={gray!18},
]
% p = 0.5 : divergent, partial sums of 1/sqrt(n)
\addplot[trapred, thick, mark=*, mark size=1.1pt]
    coordinates {(1,1)(2,1.7071)(3,2.2845)(4,2.7845)(5,3.2317)(6,3.6400)};
\addlegendentry{$p=\tfrac12$ (diverges, climbs without bound)}
% p = 1 : harmonic, divergent but slow
\addplot[quotientorange, thick, mark=*, mark size=1.1pt]
    coordinates {(1,1)(2,1.5)(4,2.0833)(6,2.45)(8,2.7179)(10,2.9290)(14,3.2516)(18,3.4951)(20,3.5977)};
\addlegendentry{$p=1$ (harmonic, diverges slowly)}
% p = 2 : convergent, approaches pi^2/6 = 1.6449
\addplot[curveblue, thick, mark=*, mark size=1.1pt]
    coordinates {(1,1)(2,1.25)(3,1.3611)(4,1.4236)(6,1.4914)(8,1.5274)(10,1.5498)(14,1.5760)(20,1.5962)};
\addlegendentry{$p=2$ (converges to $\pi^2/6\approx1.645$)}
\addplot[curveblue, dashed, domain=0:20.5] {1.6449};
\end{axis}
\end{tikzpicture}
\caption{Partial sums of three $p$-series. For $p=2$ they flatten toward a ceiling; for $p=1$ and $p=\tfrac12$ they keep rising, the half-power case visibly faster. Convergence is decided entirely by whether $p$ clears $1$.}
\label{ch29:fig-pseries}
\end{figure}

\begin{connectionbox}
The convergent $p$-series secretly compute famous constants. Euler stunned Europe in 1735 by proving $\sum_{n=1}^\infty \tfrac{1}{n^2} = \tfrac{\pi^2}{6}$ -- a fact the integral test alone cannot deliver (it tells you \emph{that} the sum is finite, not \emph{what} it equals). As a function of $p$, this object is the Riemann zeta function $\zeta(p) = \sum n^{-p}$, whose behavior encodes the distribution of the prime numbers. Convergence is only the doorway; behind it lies some of the deepest mathematics there is.
\end{connectionbox}

\section{Comparison Tests: Lean on a Series You Already Understand}
\label{ch29:comparison}

We now have two reliable yardsticks -- geometric series and $p$-series. The comparison tests let us measure unfamiliar series against these known ones. The logic is exactly what your intuition expects.

\begin{intuitionbox}
If your terms are smaller (term by term) than the terms of a series you already know converges, your sum is trapped underneath a finite total -- it has nowhere to go but to a finite value. And if your terms are bigger than those of a series you know diverges, your sum is pushed up past every bound. Smaller than finite is finite; bigger than infinite is infinite. That's the whole idea.
\end{intuitionbox}

\begin{theorem}[Direct comparison test]
\label{ch29:thm-comparison}
Suppose $0 \le a_n \le b_n$ for all $n$ (beyond some starting point).
\begin{itemize}
\item If $\sum b_n$ converges, then $\sum a_n$ converges.
\item If $\sum a_n$ diverges, then $\sum b_n$ diverges.
\end{itemize}
\end{theorem}

\begin{proof}
All terms are nonnegative, so both sequences of partial sums are increasing. Let $A_n=\sum_{k\le n} a_k$ and $B_n = \sum_{k\le n} b_k$; then $A_n \le B_n$. If $\sum b_n$ converges to $B$, then $A_n \le B_n \le B$, so $(A_n)$ is increasing and bounded above, hence convergent. The second bullet is the contrapositive of the first.
\end{proof}

\begin{example}
Does $\displaystyle\sum_{n=1}^\infty \frac{1}{n^2 + 3n + 1}$ converge? The denominator is bigger than $n^2$, so each term is \emph{smaller} than $\tfrac{1}{n^2}$:
\[
0 < \frac{1}{n^2+3n+1} < \frac{1}{n^2}.
\]
Since $\sum \tfrac1{n^2}$ converges ($p=2>1$), our smaller series converges too. We never needed its value -- only that it's trapped under something finite.
\end{example}

The direct comparison demands a clean term-by-term inequality, which can be fiddly to arrange (you may have to make the smaller series smaller, or the bigger one bigger, in just the right direction). Often it is easier to compare \emph{growth rates} rather than the terms themselves. That is the limit comparison test, and it is the friendlier workhorse.

\begin{theorem}[Limit comparison test]
\label{ch29:thm-limcomparison}
Let $a_n > 0$ and $b_n > 0$, and suppose
\[
\lim_{n\to\infty} \frac{a_n}{b_n} = L, \qquad 0 < L < \infty.
\]
Then $\sum a_n$ and $\sum b_n$ either both converge or both diverge.
\end{theorem}

\begin{whybox}
Why does a finite, nonzero limit force the two series to share a fate? Because $\tfrac{a_n}{b_n}\to L$ means that eventually $a_n$ is sandwiched between $\tfrac{L}{2}b_n$ and $2L\,b_n$ -- the two series are, in the long run, just constant multiples of each other. And multiplying every term of a series by a fixed positive constant cannot change whether it converges (it only scales the total). Convergence is decided by the \emph{tail}, by how fast terms shrink, and a constant factor does not change a shrink rate. So $\sum a_n$ and $\sum b_n$ live or die together. The limit comparison test is really the statement ``convergence cares about growth rate, not constant multiples'' made operational. You don't have to win a delicate inequality; you only have to identify which known series your terms \emph{grow like}.
\end{whybox}

\begin{focusbox}
\textbf{How to pick $b_n$ in practice.} Strip your term down to its dominant behavior for large $n$ -- keep only the fastest-growing parts of numerator and denominator. That skeleton is your comparison series. Example: $a_n = \dfrac{3n^2 - n}{4n^4 + 5}$ behaves like $\dfrac{3n^2}{4n^4} = \dfrac{3}{4n^2}$, so compare with $b_n = \tfrac{1}{n^2}$.
\end{focusbox}

\begin{example}
Test $\displaystyle\sum_{n=1}^\infty \frac{3n^2 - n}{4n^4 + 5}$. The dominant behavior is $\tfrac{1}{n^2}$, so take $b_n = \tfrac1{n^2}$:
\[
\frac{a_n}{b_n} = \frac{3n^2 - n}{4n^4 + 5}\cdot n^2 = \frac{3n^4 - n^3}{4n^4 + 5} \longrightarrow \frac{3}{4}.
\]
The limit $\tfrac34$ is finite and nonzero, and $\sum\tfrac1{n^2}$ converges, so our series converges. Notice how much cleaner this is than fighting for a direct inequality.
\end{example}

\begin{example}[A divergence by limit comparison]
Test $\displaystyle\sum_{n=1}^\infty \frac{1}{\sqrt{n^2+1}}$. For large $n$, $\sqrt{n^2+1}\approx n$, so compare with $b_n=\tfrac1n$:
\[
\frac{a_n}{b_n} = \frac{n}{\sqrt{n^2+1}} = \frac{1}{\sqrt{1+1/n^2}} \longrightarrow 1.
\]
Finite, nonzero. Since $\sum\tfrac1n$ diverges (harmonic), so does our series. The square root barely changed anything; the term still shrinks ``like $1/n$,'' and that is fatal.
\end{example}

\section{The Ratio and Root Tests: When a Series Is Secretly Geometric}
\label{ch29:ratio-root}

The geometric series converges precisely because the ratio of consecutive terms is a constant $r$ with $\abs r < 1$. The \textbf{ratio test} takes this idea and applies it to series that are only \emph{eventually}, \emph{approximately} geometric -- which, it turns out, is an enormous and important class, including every series built from factorials and exponentials.

\begin{intuitionbox}
Look at the ratio $\tfrac{a_{n+1}}{a_n}$ of one term to the previous. For a true geometric series this ratio is always $r$. For a general series it varies, but if it \emph{settles} toward some limit $L$, then far out in the tail the series behaves like a geometric series with ratio $L$. So if $L<1$, the tail decays geometrically and the series converges; if $L>1$, the terms eventually \emph{grow}, dooming the sum.
\end{intuitionbox}

\begin{theorem}[Ratio test]
\label{ch29:thm-ratio}
For a series $\sum a_n$ with $a_n \neq 0$, suppose $\displaystyle\lim_{n\to\infty}\left|\frac{a_{n+1}}{a_n}\right| = L$.
\begin{itemize}
\item If $L < 1$, the series converges (absolutely).
\item If $L > 1$ (or $L=\infty$), the series diverges.
\item If $L = 1$, the test is inconclusive -- it tells you nothing.
\end{itemize}
\end{theorem}

\begin{proof}[Why it works]
Suppose $L<1$. Pick any $r$ with $L < r < 1$. Since the ratios approach $L$, beyond some index $N$ we have $\abs{a_{n+1}} < r\,\abs{a_n}$. Then by induction $\abs{a_{N+k}} < r^k \abs{a_N}$, so the tail of the series is dominated, term by term, by the geometric series $\abs{a_N}\sum r^k$, which converges because $r<1$. By direct comparison the tail converges, hence the whole series converges absolutely. If instead $L>1$, the terms eventually grow, so $a_n \not\to 0$ and the $n$th-term test kills it. When $L=1$, the comparison geometric series sits exactly on its own boundary and tells us nothing.
\end{proof}

\begin{example}[Factorials are the ratio test's home turf]
Test $\displaystyle\sum_{n=1}^\infty \frac{2^n}{n!}$. Form the ratio and watch the factorial demolish the exponential:
\[
\left|\frac{a_{n+1}}{a_n}\right| = \frac{2^{n+1}}{(n+1)!}\cdot\frac{n!}{2^n} = \frac{2}{n+1} \longrightarrow 0.
\]
$L = 0 < 1$, so the series converges -- and fast. (Its sum, foreshadowing the next chapter, is $e^2 - 1$.) Whenever you see factorials, reach for the ratio test first; the cancellation is almost always this clean. Figure~\ref{ch29:fig-ratio} shows the terms diving under a geometric envelope, which is the whole mechanism behind the test.
\end{example}

\begin{example}[Why $L=1$ is genuinely useless]
Apply the ratio test to $\sum \tfrac1n$ (diverges) and to $\sum \tfrac1{n^2}$ (converges). Both give
\[
\sum\tfrac1n:\ \frac{1/(n+1)}{1/n}=\frac{n}{n+1}\to 1, \qquad
\sum\tfrac1{n^2}:\ \frac{n^2}{(n+1)^2}\to 1.
\]
The ratio test returns $L=1$ for both, yet one converges and one does not. This is concrete proof that $L=1$ carries no information -- you must switch to another test (here, the $p$-series rule). The ratio test is blind precisely at the borderline of polynomial decay.
\end{example}

The \textbf{root test} is the ratio test's close cousin, better suited when the $n$th term is an $n$th power. Instead of comparing neighbors, it extracts the ``geometric ratio'' directly by taking an $n$th root.

\begin{theorem}[Root test]
\label{ch29:thm-root}
For $\sum a_n$, suppose $\displaystyle\lim_{n\to\infty}\sqrt[n]{\abs{a_n}} = L$. If $L<1$ the series converges absolutely; if $L>1$ (or $\infty$) it diverges; if $L=1$ it is inconclusive.
\end{theorem}

\begin{whybox}
The root test works for the same reason as the ratio test, even more directly. If $\sqrt[n]{\abs{a_n}}\to L$ then for large $n$, $\abs{a_n}\approx L^n$ -- the term \emph{is} essentially the $n$th term of a geometric series with ratio $L$. So $L<1$ means geometric decay and convergence; $L>1$ means geometric growth and divergence. The root test cuts out the middleman: rather than inferring the geometric ratio from consecutive terms, it reads it straight off the size of $a_n$ itself. Use it whenever $a_n$ wears an exponent of $n$, e.g. $a_n = \big(\tfrac{n}{2n+1}\big)^n$.
\end{whybox}

\begin{example}
Test $\displaystyle\sum_{n=1}^\infty \left(\frac{n}{2n+1}\right)^n$. The $n$th power screams ``root test'':
\[
\sqrt[n]{\abs{a_n}} = \frac{n}{2n+1} \longrightarrow \frac12 < 1,
\]
so the series converges. A ratio-test attack here would be a mess; the root test reads the answer off in one line.
\end{example}

\begin{figure}[H]
\centering
\begin{tikzpicture}
\begin{axis}[
    width=12cm, height=6cm,
    xlabel={$n$}, ylabel={term size $|a_n|$},
    xmin=0.5, xmax=10.5, ymin=0, ymax=0.85,
    xtick={1,2,3,4,5,6,7,8,9,10},
    samples=60, domain=1:10.5,
    legend pos=north east, legend cell align=left,
    grid=both, grid style={gray!18},
]
% terms of sum 2^n/n!  : a_n = 2^n/n!
\addplot[only marks, mark=*, mark size=1.5pt, curveblue]
    coordinates {(1,2)(2,2)(3,1.3333)(4,0.6667)(5,0.2667)(6,0.0889)(7,0.0254)(8,0.00635)(9,0.001411)(10,0.000282)};
\addlegendentry{terms $2^n/n!$}
% geometric envelope past n=N: here once ratio < r=0.5, terms < a_5*(0.5)^{n-5}
\addplot[tangentred, dashed, very thick, domain=4:10.5] {0.6667*pow(0.5,x-4)};
\addlegendentry{geometric envelope $r=\tfrac12$ past $n=4$}
\end{axis}
\end{tikzpicture}
\caption{The ratio test in action on $\sum 2^n/n!$. The ratio $2/(n+1)$ drops below $\tfrac12$ once $n\ge 4$, so from there on every term is bounded by a decaying geometric series (dashed). Being trapped under a convergent geometric tail is exactly what forces convergence.}
\label{ch29:fig-ratio}
\end{figure}

\section{Alternating Series and the Sign That Saves You}
\label{ch29:alternating}

So far our terms have mostly been positive. Something remarkable happens when the signs alternate $+,-,+,-,\dots$: the back-and-forth lets partial cancellation rescue a series that would diverge if all its terms were positive. The exemplar is the \textbf{alternating harmonic series}:
\[
1 - \frac{1}{2} + \frac{1}{3} - \frac{1}{4} + \frac{1}{5} - \cdots = \sum_{n=1}^\infty \frac{(-1)^{n-1}}{n}.
\]
The positive version (the plain harmonic series) diverges to infinity. But with alternating signs, this one \emph{converges} -- to $\ln 2$, in fact. The minus signs claw back the runaway growth.

\begin{intuitionbox}
Watch the partial sums of an alternating series whose terms shrink. You go up by $1$, down by $\tfrac12$, up by $\tfrac13$, down by $\tfrac14$. Each step overshoots the eventual limit, but by less than the last step (because the steps are shrinking). The partial sums zig-zag with ever-tightening swings -- like a pendulum losing energy, homing in on a single value squeezed between every up-step and every down-step.
\end{intuitionbox}

\begin{figure}[H]
\centering
\begin{tikzpicture}
\begin{axis}[
    width=12cm, height=6cm,
    xlabel={$n$}, ylabel={partial sum $S_n$},
    xmin=0.5, xmax=8.5, ymin=0.4, ymax=1.1,
    xtick={1,2,3,4,5,6,7,8},
    legend pos=north east, legend cell align=left,
    grid=both, grid style={gray!18},
]
\addplot[curveblue, thick, mark=*, mark size=1.5pt]
    coordinates {(1,1)(2,0.5)(3,0.8333)(4,0.5833)(5,0.7833)(6,0.6167)(7,0.7595)(8,0.6345)};
\addlegendentry{$S_n$ (alternating harmonic)}
\addplot[tangentred, dashed, thick, domain=0.5:8.5] {0.6931};
\addlegendentry{limit $\ln 2 \approx 0.693$}
\end{axis}
\end{tikzpicture}
\caption{Partial sums of $1 - \tfrac12 + \tfrac13 - \cdots$ oscillate with shrinking amplitude around $\ln 2$. The odd partial sums approach from above, the even ones from below, and they trap the limit between them.}
\label{ch29:fig-alternating}
\end{figure}

\begin{theorem}[Alternating series test, Leibniz]
\label{ch29:thm-alternating}
If the positive numbers $b_n$ satisfy (i) $b_{n+1} \le b_n$ (decreasing) and (ii) $b_n \to 0$, then the alternating series $\sum_{n=1}^\infty (-1)^{n-1} b_n$ converges. Moreover, the error from stopping at the $n$th term is no larger than the first omitted term: $\abs{S - S_n} \le b_{n+1}$.
\end{theorem}

\begin{proof}
Consider the even partial sums. Grouping in pairs, $S_{2n} = (b_1 - b_2) + (b_3 - b_4) + \cdots + (b_{2n-1}-b_{2n})$; each pair is $\ge 0$ since $b_k$ decreases, so $(S_{2n})$ is increasing. Regroup the other way, $S_{2n} = b_1 - (b_2 - b_3) - \cdots - b_{2n}$; each subtracted pair is $\ge 0$, so $S_{2n} \le b_1$. An increasing sequence bounded above converges -- call the limit $S$. Finally $S_{2n+1} = S_{2n} + b_{2n+1}$, and since $b_{2n+1}\to 0$ the odd sums approach the same $S$. So all partial sums converge to $S$. The error bound follows because $S$ always lies between consecutive partial sums, which differ by exactly $b_{n+1}$.
\end{proof}

\begin{aestheticbox}
That error bound is a small marvel. It says: to approximate the sum to within $0.001$, just keep adding terms until the terms themselves drop below $0.001$, then stop -- the leftover error is guaranteed smaller than the first term you threw away. No other elementary test hands you such a clean, automatic accuracy guarantee. This is why alternating series are beloved in numerical work: they come with their own honesty built in.
\end{aestheticbox}

\section{Absolute vs.\ Conditional Convergence: A Subtle Hierarchy}
\label{ch29:absolute}

The alternating harmonic series converges, but only thanks to its signs -- strip them away and it diverges. Contrast that with $\sum \tfrac{(-1)^{n-1}}{n^2}$, which converges even with all signs made positive (it's the $p$-series $p=2$). These two convergences are not the same \emph{kind}, and the difference matters enormously.

\begin{definition}[Absolute and conditional convergence]
\label{ch29:def-absolute}
A series $\sum a_n$ \textbf{converges absolutely} if $\sum \abs{a_n}$ converges. It \textbf{converges conditionally} if $\sum a_n$ converges but $\sum \abs{a_n}$ diverges -- that is, the convergence depends on cancellation among the signs.
\end{definition}

\begin{theorem}[Absolute convergence implies convergence]
\label{ch29:thm-abs-implies}
If $\sum \abs{a_n}$ converges, then $\sum a_n$ converges.
\end{theorem}

\begin{proof}
The clean trick: $0 \le a_n + \abs{a_n} \le 2\abs{a_n}$ (the middle is either $0$ or $2\abs{a_n}$). Since $\sum 2\abs{a_n}$ converges, by direct comparison $\sum (a_n + \abs{a_n})$ converges. Then $\sum a_n = \sum (a_n + \abs{a_n}) - \sum \abs{a_n}$ is a difference of two convergent series, hence converges.
\end{proof}

\begin{focusbox}
This theorem is what makes the ratio and root tests legitimate. When they report ``$L<1$,'' they are really proving $\sum\abs{a_n}$ converges -- absolute convergence -- which by the theorem guarantees ordinary convergence. So those tests handle signs for free: you can throw a sign-changing series at them and they just work, because absolute convergence is the stronger, sign-blind notion.
\end{focusbox}

\begin{happeningbox}
What is the \emph{real} difference between the two kinds? An absolutely convergent series converges because its terms are genuinely, sustainably small -- it would converge even if you erased every minus sign. Its sum is robust: you may reorder, regroup, and rearrange the terms however you like and always reach the same total, exactly as with a finite sum. A conditionally convergent series is a far stranger animal. It converges only because of a delicate \emph{balancing act} between positive and negative terms; the positives alone sum to $+\infty$ and the negatives alone to $-\infty$, and the finite total is just the tug-of-war between two infinities. Disturb that balance -- reorder the terms -- and you can make the sum come out to \emph{any number you wish}. This is Riemann's astonishing Rearrangement Theorem. For the alternating harmonic series, one rearrangement gives $\ln 2$, another gives $5$, another gives $-\pi$, all from the very same terms in a different order. Conditional convergence is convergence held together by string and good timing; absolute convergence is convergence built of stone.
\end{happeningbox}

\begin{historybox}
Bernhard Riemann discovered this rearrangement phenomenon in the 1850s, and it shocked a generation that had assumed infinite sums obeyed the same rules as finite ones. They do not. Commutativity of addition -- $a+b=b+a$, extended to ``order never matters'' -- silently fails for conditionally convergent series. This was one of several nineteenth-century jolts (alongside the divergence surprises and pathological functions) that convinced mathematicians infinite processes needed the rigorous foundation that Cauchy and Weierstrass were busy building. The partial-sum definition we opened with is the inheritance of exactly that crisis: a way to reason about the infinite without being ambushed by it.
\end{historybox}

\section{Choosing a Test: A Strategy, Not a Lottery}
\label{ch29:strategy}

With seven or eight tests in hand, the real skill is knowing which to reach for. The choice is not arbitrary -- the \emph{form} of the term usually tells you. Here is the decision flow an experienced eye runs, almost unconsciously.

\begin{focusbox}
\textbf{A practical order of attack.}
\begin{enumerate}
\item \textbf{Glance at the terms.} Do they fail to go to $0$? Then it diverges by the $n$th-term test (Sec.~\ref{ch29:nth-term}). One look -- always do this first.
\item \textbf{Recognize the species.} Is it geometric ($a r^n$) or telescoping (a difference $b_k - b_{k+1}$)? Then sum it \emph{exactly} (Secs.~\ref{ch29:geometric},~\ref{ch29:telescoping}).
\item \textbf{Spot factorials or $n$th powers.} Factorials or exponentials $\Rightarrow$ ratio test. A whole term raised to the $n$th power $\Rightarrow$ root test (Sec.~\ref{ch29:ratio-root}).
\item \textbf{See a rational / algebraic term} (powers of $n$, roots): find its dominant $p$-power and use limit comparison against a $p$-series (Secs.~\ref{ch29:comparison},~\ref{ch29:integral-test}).
\item \textbf{Notice alternating signs} $(-1)^n$: check absolute convergence first (it's stronger and easier); if that fails, try the alternating series test (Secs.~\ref{ch29:alternating},~\ref{ch29:absolute}).
\item \textbf{If $f(n)$ has a clean antiderivative} ($\tfrac1{n\ln n}$, $n e^{-n^2}$): the integral test may be the most direct route (Sec.~\ref{ch29:integral-test}).
\end{enumerate}
\end{focusbox}

\begin{trapbox}
The classic blunders, collected so you can avoid them:
\begin{itemize}
\item Concluding \emph{convergence} from $a_n\to 0$. The $n$th-term test only proves divergence; it is silent on convergence. (Harmonic series.)
\item Declaring victory when the ratio or root test returns $L=1$. That is the \emph{inconclusive} verdict -- switch tests, do not guess.
\item Using the integral test when $f$ is not eventually decreasing. The hypothesis is essential; without it the rectangle-curve sandwich falls apart.
\item Using limit comparison with $L=0$ or $L=\infty$ carelessly. The clean ``share a fate'' conclusion needs $0<L<\infty$; the boundary cases give only one-directional information.
\item Reordering a conditionally convergent series as if order didn't matter. It does (Riemann).
\end{itemize}
\end{trapbox}

\begin{labbox}
Classify each by hand, naming the test and the reasoning before you compute. (a) $\sum \dfrac{n}{n^2+1}$. \;(b) $\sum \dfrac{n!}{n^n}$. \;(c) $\sum \dfrac{(-1)^n}{\ln n}$ (from $n=2$). \;(d) $\sum \dfrac{\cos n}{n^2}$.

\emph{Hints to check yourself:} (a) limit-compare with $\tfrac1n$, diverges. (b) ratio test, $\big(\tfrac{n}{n+1}\big)^n \to \tfrac1e<1$, converges. (c) terms decrease to $0$, alternating series test, converges -- but $\sum\tfrac{1}{\ln n}$ diverges, so only \emph{conditionally}. (d) $\big|\tfrac{\cos n}{n^2}\big|\le \tfrac{1}{n^2}$, comparison gives \emph{absolute} convergence.
\end{labbox}

\begin{applicationbox}
Series are not museum pieces. In physics, the potential energy of a crystal lattice is an infinite (often conditionally convergent!) sum over ions, and getting the rearrangement right is a real computational concern. In probability, the expected value of a random variable is a series $\sum x_k p_k$, and whether it converges \emph{absolutely} decides whether the expectation is even well-defined. In machine learning, a discounted reward $\sum_{t=0}^\infty \gamma^t r_t$ is a geometric-type series whose convergence (requiring $\abs\gamma<1$) is what makes reinforcement learning's value functions finite. And every special function -- Bessel functions, the error function, $\zeta(s)$ -- is delivered to your computer as a series whose convergence rate determines how many terms it must add to hit the precision you asked for. The tests in this chapter are, quietly, what make numerical computation possible.
\end{applicationbox}

\section{Exercises}
\label{ch29:exercises}

\begin{exercise}[Sum it exactly]
Evaluate each convergent series in closed form.
(a) $\displaystyle\sum_{n=0}^\infty \frac{3}{4^n}$. \quad
(b) $\displaystyle\sum_{n=1}^\infty \left(\frac{2}{3}\right)^{n}$. \quad
(c) $\displaystyle\sum_{n=1}^\infty \frac{1}{(2n-1)(2n+1)}$ (telescope via $\tfrac12\big(\tfrac{1}{2n-1}-\tfrac{1}{2n+1}\big)$). \quad
(d) Express $0.1\overline{6}$ (that is, $0.16666\ldots$) as a fraction using a geometric series.
\end{exercise}

\begin{exercise}[The $n$th-term test, used correctly]
Determine which of these diverge by the $n$th-term test, and explain why the test is \emph{silent} on the others.
(a) $\sum \dfrac{n^2}{n^2+5}$. \;
(b) $\sum \dfrac{1}{n^2+5}$. \;
(c) $\sum \cos\!\big(\tfrac1n\big)$. \;
(d) $\sum (-1)^n \dfrac{n}{n+1}$.
\end{exercise}

\begin{exercise}[$p$-series and comparison]
Decide convergence, naming your comparison series each time.
(a) $\sum \dfrac{1}{n^{3/2}}$. \;
(b) $\sum \dfrac{1}{\sqrt[3]{n}}$. \;
(c) $\sum \dfrac{n+4}{n^3 - 2n + 7}$. \;
(d) $\sum \dfrac{1}{\sqrt{n^3+1}}$.
\end{exercise}

\begin{exercise}[Ratio and root]
Apply the ratio or root test, whichever is cleaner, and say what $L$ is.
(a) $\sum \dfrac{5^n}{n!}$. \;
(b) $\sum \dfrac{n!}{10^n}$. \;
(c) $\sum \dfrac{n^2}{2^n}$. \;
(d) $\sum \left(\dfrac{3n+1}{4n+5}\right)^{n}$.
\end{exercise}

\begin{exercise}[Integral test]
Use the integral test to settle each, and verify the hypotheses (positive, decreasing) hold.
(a) $\displaystyle\sum_{n=2}^\infty \frac{1}{n\ln n}$. \;
(b) $\displaystyle\sum_{n=2}^\infty \frac{1}{n(\ln n)^2}$. \;
(c) $\displaystyle\sum_{n=1}^\infty n\,e^{-n^2}$.
Comment on how (a) and (b) differ despite looking nearly identical -- where is the borderline now?
\end{exercise}

\begin{exercise}[Alternating, absolute, conditional]
For each, decide whether it converges absolutely, converges conditionally, or diverges. Justify all three possibilities you rule out or confirm.
(a) $\displaystyle\sum_{n=1}^\infty \frac{(-1)^{n-1}}{\sqrt n}$. \;
(b) $\displaystyle\sum_{n=1}^\infty \frac{(-1)^{n-1}}{n^2}$. \;
(c) $\displaystyle\sum_{n=1}^\infty \frac{(-1)^{n-1}\,n}{n+1}$. \;
(d) $\displaystyle\sum_{n=2}^\infty \frac{(-1)^n}{n\ln n}$.
\end{exercise}

\begin{exercise}[Error estimate]
For the alternating series $\displaystyle\sum_{n=1}^\infty \frac{(-1)^{n-1}}{n^3}$, how many terms must you add to guarantee the partial sum is within $10^{-3}$ of the true value? Use the alternating-series error bound, and explain why that bound applies here.
\end{exercise}

\begin{exercise}[Conceptual]
Answer in a sentence or two each, with reasons.
(a) True or false: if $a_n \to 0$ then $\sum a_n$ converges. Give the canonical counterexample.
(b) True or false: if $\sum a_n$ converges then $a_n \to 0$. Which theorem is this?
(c) Can a series of positive terms converge conditionally? Why or why not?
(d) If $\sum a_n$ and $\sum b_n$ both diverge, must $\sum(a_n+b_n)$ diverge? Find an example or prove it.
\end{exercise}

\begin{exercise}[Prove it]
Prove that if $\sum a_n$ converges and $\sum b_n$ diverges, then $\sum (a_n + b_n)$ diverges. (Hint: argue by contradiction using the fact that the difference of two convergent series converges.)
\end{exercise}

\begin{exercise}[Prove it -- the heart of the matter]
Prove the $n$th-term test from scratch: if $\sum_{k=1}^\infty a_k = S$ (converges), then $a_n \to 0$. Your proof should use only the definition of convergence via partial sums (Definition~\ref{ch29:def-convergence}) and the fact that if $S_n \to S$ then $S_{n-1}\to S$. Then explain in one paragraph, in your own words, why the converse fails, citing a specific series.
\end{exercise}

\begin{exercise}[Geometric in disguise]
A patient takes a $100$~mg dose of a drug each day. Each day, the body eliminates $25\%$ of whatever drug is present, so $75\%$ carries over to the next day. Model the long-run amount of drug in the body just after a dose as an infinite series, identify it as geometric, and find the steady-state amount. (This is exactly the reinforcement-learning discounted-sum structure from the application box, wearing a lab coat.)
\end{exercise}

\begin{exercise}[Build the borderline]
We saw $\sum \tfrac1n$ diverges and $\sum\tfrac1{n^2}$ converges. Investigate the finer scale between them:
(a) Argue that $\sum \tfrac{1}{n^{1.0001}}$ converges and $\sum\tfrac{1}{n^{0.9999}}$ diverges -- the borderline at $p=1$ is sharp.
(b) Using the integral test, show $\sum \tfrac{1}{n\ln n}$ still diverges, so $\tfrac1n$ is not even the ``last'' divergent term -- there is a whole ladder of ever-slower divergences below it. Reflect briefly on what this says about the idea of a single ``boundary'' between convergence and divergence.
\end{exercise}
